\documentclass[12pt,a4paper]{article}
\usepackage[utf8]{inputenc}
\usepackage{geometry}
\geometry{margin=1in}
\usepackage{hyperref}
\usepackage{graphicx}
\usepackage{booktabs}
\usepackage{longtable}

\title{\textbf{Landslide Identification Using Machine Learning\\ \large Project Documentation}}
\author{}
\date{\today}

\begin{document}

\maketitle
\tableofcontents
\newpage

\section{Project Overview}
This project aims to automate the detection of landslides using satellite imagery and historical event data. Landslides are devastating natural disasters, and quickly identifying them helps direct emergency operations and damage assessment. This system operates in two phases:
\begin{enumerate}
    \item \textbf{Phase 1 (Image Classification):} A Convolutional Neural Network (CNN) looks at a satellite image and decides if there is a landslide in it.
    \item \textbf{Phase 2 (HMM Future Predictor):} If a landslide is detected, a Hidden Markov Model (HMM) analyzes historical data for that specific country to determine the landslide type and predict the probability of future landslide events.
\end{enumerate}

\section{Problem Statement}
Every year, landslides happen in mountainous and rain-heavy areas, causing loss of life and property. Currently, experts manually inspect satellite images to locate landslides. This process is slow, expensive, and scales poorly in large regions.

The goal of this project is to create an automated, two-phase machine learning pipeline that not only detects current landslides from satellite images quickly, but also acts as an early warning system by forecasting what type of landslide might naturally follow in that geological region based on historical records.

\section{Dataset Explanation}
The system uses two different sets of data to learn:
\begin{itemize}
    \item \textbf{HR-GLDD (High-Resolution Global Landslide Detector Dataset):} This data comes from Zenodo. It provides high-resolution satellite image patches. During preprocessing, images with significant landslide presence were marked as "landslide", and patches with no landslides were marked as "non\_landslide" (Total: 3,439 images covering train, validation, and test splits).
    \item \textbf{NASA Global Landslide Catalog:} Used for Phase 2, this is a record of 9,471 real landslide events over 28 years. It includes the date, country, the type of landslide (e.g., Mudslide, Rockfall), and what triggered it (e.g., Heavy Rain, Earthquake).
\end{itemize}

\section{Machine Learning Concepts Used}
To understand the project, you should know these fundamental terms:
\begin{itemize}
    \item \textbf{Machine Learning (ML):} Teaching a computer to learn patterns from examples instead of giving it explicit rules.
    \item \textbf{Dataset:} A collection of examples (like our satellite images) used to teach the computer.
    \item \textbf{Feature:} A specific detail in the data. In images, "features" are textures, colors, or shapes that the model learns to associate with a landslide.
    \item \textbf{Training:} The phase where the model studies the dataset, makes guesses, and corrects its own mistakes to get better.
    \item \textbf{Testing/Validation:} Giving the model fresh data it hasn't seen before to check if it truly learned the task.
    \item \textbf{Model:} The mathematical engine resulting from training. When we input a new image, the model gives us an answer.
    \item \textbf{Algorithm:} The mathematical steps the computer follows to learn. We use Convolutional Neural Networks (CNNs) for images, and Hidden Markov Models (HMM) for sequences of events.
    \item \textbf{Accuracy:} The percentage of correct predictions the model makes.
    \item \textbf{Loss:} A score that measures how wrong the model is while training. The goal is to minimize this score.
\end{itemize}

\section{Explanation of Each File/Folder in the Project}
The project code is structured functionally:
\begin{itemize}
    \item \texttt{project/config.py}: Holds the master settings for the project (like file paths, number of training cycles, image sizes).
    \item \texttt{project/phase1\_alexnet/}: Contains the code for the Image Looker.
    \begin{itemize}
        \item \texttt{dataset.py}: Prepares and loads the images for training.
        \item \texttt{model.py}: Defines the shapes of the Neural Networks (AlexNet and EfficientNet-B3).
        \item \texttt{train.py / evaluate.py / predict.py / ensemble\_predict.py}: Scripts to teach the model, test its accuracy, and run predictions using combinations of models.
    \end{itemize}
    \item \texttt{project/phase2\_hmm/}: Contains the code for the Future Predictor.
    \begin{itemize}
        \item \texttt{data\_preprocessing.py}: Cleans the NASA data.
        \item \texttt{hmm\_model.py / hmm\_predict.py}: Teaches the HMM using the NASA history and generates future forecasts.
    \end{itemize}
    \item \texttt{project/pipeline/run\_pipeline.py}: The main traffic controller. It feeds the image to Phase 1, and if Phase 1 says "Landslide", it passes the baton to Phase 2 to predict the future.
    \item \texttt{presentation.md / content.md}: Contains detailed records of the project's logic, steps taken, and past results.
\end{itemize}

\section{Model Architecture / Workflow}
\textbf{Phase 1 Architecture: Ensemble Deep Learning}\\
To detect landslides, the project uses an "Ensemble". It combines two "brains" working at the same time:
1. \textbf{EfficientNet-B3 (60\% vote):} A modern, efficient network excellent at noticing subtle details (high recall).
2. \textbf{AlexNet (40\% vote):} A classic network that helps double-check the results and prevents false alarms (high precision).
When a new image is provided, both brains analyze it. Their probabilities are mathematically blended. If the blended confidence is above 46.7\%, it triggers Phase 2.

\textbf{Phase 2 Architecture: Hidden Markov Model (HMM)}\\
When activated, the HMM fetches the country's landslide history. It uses 8 "hidden states" (representing underlying geological regimes) to predict the most likely type of landslide occurring now, and forecasts the probability of future landslides up to three steps ahead.

\section{Training Process}
The CNN (Phase 1) was trained on 2,190 satellite images using a technique called \textbf{Transfer Learning}. We didn't start from scratch; we used a model that already studied 1.2 million internet images, allowing it to learn the landslide shapes much faster. We trained it over 40 ``epochs" (passes over the data).

The HMM (Phase 2) was trained using the \textbf{Baum-Welch algorithm} for 100 iterations. It scanned through the sequences of 9,471 historical NASA events to learn which landslides tend to follow each other.

\section{Evaluation Metrics}
The models are judged using the following metrics:
\begin{itemize}
    \item \textbf{Accuracy:} Overall correctness across all predictions.
    \item \textbf{Precision:} When the model says "Landslide," how often is it right?
    \item \textbf{Recall:} Out of all the real landslides, how many did the model successfully find?
    \item \textbf{F1-Score:} A balanced average between Precision and Recall.
    \item \textbf{ROC-AUC:} Measures the model's ability to distinguish between classes across all possible thresholds.
\end{itemize}

\section{Results Overview}
The project improved dramatically over 5 major versions, culminating in the final ensemble:
\begin{table}[h]
\centering
\begin{tabular}{@{}llcc@{}}
\toprule
\textbf{Version} & \textbf{Model / Method} & \textbf{Improvement} & \textbf{Accuracy} \\ \midrule
Result 1 & AlexNet (Scratch) & Baseline & 78.54\% \\
Result 2 & AlexNet (Pretrained) & Transfer Learning & 80.11\% \\
Result 3 & EfficientNet-B3 & Better Architecture & 79.97\% \\
Result 4 & EfficientNet + HMM Upgrade & Confidence Calibration & 79.97\% \\
Result 5 & Ensemble + Optimal Threshold & Teamwork cancels mistakes & \textbf{82.40\%} \\ \bottomrule
\end{tabular}
\end{table}

The final ensemble boasts a recall of 78.67\% and an F1-Score of 72.97\%, proving it correctly identifies the vast majority of landslides while keeping false alarms to a minimum.

\end{document}
